\documentclass{situationist-preprint}

\title{Cantor's Problem and Its Status}
\msid{MTH-06-c}
\msfamily{MTH-06}
\msdiscipline{Philosophy of Mathematics}
\mstype{Consequence}
\mskeywords{continuum hypothesis; Cantor; Hilbert's first problem; independence; open problems; multiverse}

\begin{document}
\maketitle

\begin{abstract}
Cantor's continuum problem is often described as solved by the independence results of G\"odel and Cohen, or as shown to be meaningless by them. We derive the status the problem has if, as MTH-06 holds, CH is false given new axioms that are independently motivated. The problem is neither solved in the sense of a proof from accepted axioms nor dissolved. It is answered: its answer is known on grounds of the same kind as support the axioms of set theory themselves. What remains open is a different question, the value of the continuum.
\end{abstract}

\section{Three statuses}\label{s:status}

\citet{Cantor1878} conjectured that every infinite set of reals is either countable or of the cardinality of the continuum. \citet{Hilbert1900} placed the question first among his problems. G\"odel showed that CH is consistent with ZFC, and \citet{Cohen1963} that its negation is. Three statuses have since been ascribed to the problem. It is \emph{solved} if CH or its negation is proved from accepted axioms; it is \emph{dissolved} if it has no determinate answer, as on the multiverse view of \citet{Hamkins2012}; it is \emph{answered} if there is a determinate answer that is justified, though not by proof from the axioms accepted when the problem was posed.

\begin{proposition}\label{p:status}
Given MTH-06, Cantor's problem is answered, and is neither solved nor dissolved.
\end{proposition}

\begin{proof}
It is not solved, since ZFC proves neither CH nor its negation. It is not dissolved, since by Principle P1 of MTH-06-a CH has a determinate truth value. It is answered: by Theorem~5.1 of MTH-06-a, $\neg$CH is justified given independently motivated axioms, and by Proposition~5.4 of MTH-06-a no rival axiom that implies CH is justified.
\end{proof}

The status is not peculiar to CH. The question whether every projective set of reals is Lebesgue measurable is independent of ZFC and is answered, by the general acceptance of large cardinal axioms, in the affirmative. \citet{Godel1947} predicted that the continuum problem would be settled the same way, by new axioms whose credentials are their consequences; MTH-06 holds that the prediction has been fulfilled, though by forcing axioms rather than by the large cardinals G\"odel had in view.

\section{What remains open}\label{s:open}

The answered question is Cantor's: whether there is an intermediate size between the countable and the continuum. The value of $2^{\aleph_0}$ is a further question. Martin's Maximum and Woodin's axiom $(*)$ both give $\aleph_2$ (Lemmas~4.1 and~4.3 of MTH-06-a), and Freiling's axiom of symmetry gives only that the value exceeds $\aleph_1$ (Lemma~4.4). As MTH-06-a records, the evidence favors $\aleph_2$ without establishing it.

\begin{objection}[Answered is not a status]
A question answered only relative to axioms not everyone accepts is still open; \enquote{answered} is a polite name for \enquote{settled by one school}.
\end{objection}

\begin{reply}
Every mathematical answer is relative to axioms, and the axioms of ZFC were themselves accepted partly on extrinsic grounds, as the axiom of choice was. What distinguishes an answer from a school's preference is whether the axioms that yield it are independently motivated, which is the condition MTH-06-a imposes and MTH-06-b examines. An answer of that kind is defeasible, as is the acceptance of ZFC; it is not merely sectarian.
\end{reply}

\section{Conclusion}

Cantor's problem has an answer, and the answer is that his conjecture is false: there is a size of infinite set of reals between the countable and the continuum. The answer rests on new axioms, not on a proof in ZFC, and it leaves the value of the continuum to be determined.

\begin{thebibliography}{9}\small
\bibitem[Cantor(1878)]{Cantor1878} Cantor, G. (1878). Ein Beitrag zur Mannigfaltigkeitslehre. \emph{Journal f\"ur die reine und angewandte Mathematik}, 84, 242--258.
\bibitem[Cohen(1963)]{Cohen1963} Cohen, P.~J. (1963). The independence of the continuum hypothesis. \emph{Proceedings of the National Academy of Sciences}, 50, 1143--1148.
\bibitem[G\"odel(1947)]{Godel1947} G\"odel, K. (1947). What is Cantor's continuum problem? \emph{American Mathematical Monthly}, 54, 515--525.
\bibitem[Hamkins(2012)]{Hamkins2012} Hamkins, J.~D. (2012). The set-theoretic multiverse. \emph{The Review of Symbolic Logic}, 5, 416--449.
\bibitem[Hilbert(1900)]{Hilbert1900} Hilbert, D. (1900). Mathematische Probleme. \emph{Nachrichten von der K\"oniglichen Gesellschaft der Wissenschaften zu G\"ottingen, Mathematisch-Physikalische Klasse}, 253--297.
\end{thebibliography}

\end{document}
